\section{A restricted pgmpy integration: verify the candidate, preserve the question}
\label{exec13:pgmpy-reuse}

The proposal selects a narrow pgmpy integration for generating an adjustment
candidate from a supplied, fully observed causal DAG. The project must then
verify that candidate independently before reporting it. This choice reuses
a permissively licensed graph library while preserving a scientific boundary:
a successful method call does not prove that the supplied graph describes
the bank, that the requested adjustment identifies the intended effect, or
that a numerical estimator has valid uncertainty. The selection authorizes
development of this limited prototype path; it is not evidence that an
adapter has been implemented or executed
\cite{exec13_pgmpy_reuse_primary,exec13_pgmpy_reuse_oracle}.

The source is pinned to the repository revision beginning
\texttt{1f6296fc2606}, committed on 5 October 2026 and selected by the
cutoff-bounded history query. Its project metadata declares version 1.1.2.
Fifteen inspected source, guide, test and licensing files match the selected
Git tree, in addition to their local SHA-256 identities. Source acquisition
occurred after the 7 October cutoff. The pin establishes which implementation
is discussed; neither a late repository snapshot nor a recent commit measures
maintainer responsiveness or long-term support
\cite{exec13_pgmpy_reuse_primary}.

\subsection{A guide, a graph type and an implemented method are separate evidence}

The identification guide presents minimal, variance-optimal and all-set
adjustment routes. The pinned implementation gives a more qualified picture.
Its minimal-variance branch raises \texttt{NotImplementedError}. The minimal
branch explicitly requires a single exposure and outcome, constructs a
proper backdoor graph, then obtains a minimal separator. The all-set branch
enumerates subsets of an ancestor candidate set after removing specified
causal-path and latent variables. These operations cannot be combined into
a claim that every documented graph class and every advertised adjustment
variant is implemented with the same semantics
\cite{exec13_pgmpy_reuse_primary}.

The distinction between minimal and minimum matters even without statistical
estimation. A set is inclusion-minimal when no proper subset remains valid.
It is minimum-cardinality when no valid set has fewer elements. The inspected
separator routine greedily removes dispensable members of an initial set;
that does not establish a global cardinality minimum. Nor does cardinality
alone determine estimator variance, data acquisition cost or missingness.
A front end should not translate the option name into the promise
``the cheapest or most precise set of controls''
\cite{exec13_pgmpy_reuse_primary}.

An exact graph illustrates the difference. Let $C$ cause $A$, $B$ and $Y$;
let $A$ and $B$ cause $X$; and include $X\to Y$. The two inclusion-minimal
backdoor sets are $\{C\}$ and $\{A,B\}$. Conditioning on $C$ blocks both
forks; without $C$, both $A$ and $B$ are needed. A project reference that
deletes dispensable variables from $\{A,B,C\}$ obtains either answer
depending on deletion order. That computation describes the semantics of
greedy minimality; it is not a measured pgmpy output. For a lending example,
one verified common driver may be cheaper to collect than two separate
proxies, but its actual measurement quality and availability still determine
whether the smaller set is useful \cite{exec13_pgmpy_reuse_oracle}.

The package's all-set path is also a bounded algorithmic choice. Enumerating
subsets can require work exponential in the candidate count. Restricting
the candidates to ancestors is meaningful, but a consumer must retain that
domain rather than describe the output as all imaginable observed-variable
sets. Variables irrelevant to the target, descendant restrictions and
generalized adjustment criteria introduce different enumeration questions.
The first proposed wrapper therefore excludes both all-set enumeration and
the unimplemented minimum-variance option. It asks for one candidate and
checks its stated sufficient criterion
\cite{exec13_pgmpy_reuse_primary}.

\subsection{A total-effect validator must reject the wrong question}

The inspected public validation call first checks the supported graph type
and invokes a graph-role check. The latter checks the presence of exposure
and outcome roles. The adjustment validator then tests whether parents of
the exposure remain connected to the outcome conditional on the exposure
and proposed adjustment variables. This particular call path does not
independently exclude descendants of the exposure
\cite{exec13_pgmpy_reuse_primary}.

Consider the supplied DAG $X\to M\to Y$ with $X$ a root. Its exposure
has no parents, so the inspected parent's-connection predicate is an empty
conjunction and evaluates to true, even if the proposed adjustment variable
is the mediator $M$. The role-presence check does not reject that assignment.
This is a static trace through authentic source, not an installed-engine
experiment. It suffices to show why the project cannot adopt that validation
call alone as its total-effect admission rule. The finding concerns this
specific validator and pin; it does not assert that every identification
method returns an invalid set
\cite{exec13_pgmpy_reuse_primary,exec13_pgmpy_reuse_oracle}.

The financial significance can be derived exactly. Let independent,
zero-mean, unit-variance errors generate
\[
 X=U_X,\qquad M=2X+U_M,\qquad Y=3M+U_Y.
\]
Under an intervention on $X$, the derivative of expected $Y$ is $6$.
The covariance entries are
\[
 \operatorname{Var}(X)=1,\quad \operatorname{Var}(M)=5,\quad
 \operatorname{Cov}(X,M)=2,\quad
 \operatorname{Cov}(X,Y)=6,\quad \operatorname{Cov}(M,Y)=15.
\]
The coefficient of $X$ in the population regression on $(X,M)$ is
\[
 \frac{6\cdot5-2\cdot15}{1\cdot5-2^2}=0.
\]
Conditioning on the mediator has removed the mediated path. Reporting that
coefficient as the total effect would change the causal question. If $X$
is a one-unit service intensity, $Y$ is dollars over the agreed horizon and
the common linear law is stipulated for $1{,}000$ accounts, the authored
gross response scenario is $6{,}000$ dollars, compared with zero under the
misreported coefficient. This is generated arithmetic, not bank ROI, a
binary-treatment guarantee or a measured customer response
\cite{exec13_pgmpy_reuse_oracle}.

The proposed first path consequently requires a conservative pre-action
backdoor contract. A candidate set must contain observed variables distinct
from exposure and outcome, include no descendants of the exposure, and
d-separate exposure and outcome in the graph with outgoing exposure edges
removed. This is a sufficient contract for the selected DAG setting. It
is intentionally narrower than generalized adjustment over all supported
graph classes. An otherwise useful generalized set outside the contract
receives an unsupported-scope disposition; it is not declared scientifically
impossible. The separation between a restricted sufficient criterion and
complete identification is part of the result's meaning.

Two project-owned graph calculations support the proposed verification
boundary. One tests separation after ancestral moralization. The other
enumerates simple skeleton paths, checks noncollider conditioning and checks
whether colliders have conditioned descendants. They agree on $26{,}064$
queries across all $543$ labeled DAGs on four nodes, including every ordered
exposure--outcome pair and every subset of the remaining nodes. The check
helps detect mistakes in the small reference implementation; it is not a
proof for arbitrary graph size, an upstream pgmpy test run or a statement
about statistical discovery from finite samples
\cite{exec13_pgmpy_reuse_oracle}.

\subsection{A selected DAG can conceal unresolved direction or a failed extension}

Static discovery remains a separate reuse decision. The pinned PC class
supports original, stable and parallel variants; its defaults include
parallel execution, test level $.01$, a maximum conditioning-set size of
five and a partially directed return. In the fitted method, the names
\texttt{pdag} and \texttt{cpdag} follow the same return branch.
The \texttt{dag} option calls a conversion selecting one orientation.
The option spelling alone cannot establish a complete equivalence class,
whole-graph confidence or causal identification under hidden confounding.
The source declares causal sufficiency, acyclicity and faithfulness;
particular conditional-independence tests add their own assumptions
\cite{exec13_pgmpy_reuse_primary}.

The default conditioning cap can stop search before all possible separating
sets have been considered. Test selection, data types and expert constraints
must therefore be retained alongside the graph. A per-test threshold of
$.01$ is not a $.99$ confidence level for every edge or the downstream
effect. Required and forbidden edges introduce external information;
their provenance must survive any business-language or agent-assisted
formulation step. A graph's attractive rendering cannot supply that
provenance \cite{exec13_pgmpy_reuse_primary}.

The PDAG implementation contains a particularly clear conversion boundary.
Its \texttt{to\_dag} routine warns when it cannot find a consistent
extension preserving directed edges and v-structures, then orients remaining
edges arbitrarily while avoiding cycles. The result can therefore be an
acyclic graph with different v-structures. The \texttt{to\_cpdag} route
first obtains an extension when undirected edges exist, so it cannot be
used as an independent certificate that the input represented a valid
equivalence class \cite{exec13_pgmpy_reuse_primary}.

A four-node chordless undirected cycle supplies an exact refusal fixture.
There are sixteen orientations of its four edges; two form directed cycles
and fourteen are acyclic. Every acyclic orientation has a sink with two
nonadjacent parents and therefore introduces an unshielded collider.
Consequently none preserves the undirected graph's absence of specified
colliders. The project enumerates all fourteen acyclic cases to check this
small argument. An arbitrary DAG returned after a warning would answer a
different graphical question; simply confirming acyclicity would miss the
change \cite{exec13_pgmpy_reuse_oracle}.

The initial integration therefore consumes an already supplied DAG rather
than an automatically completed discovered graph. Future discovery
integration must preserve typed directed and undirected edges, isolates,
separating sets, constraints, warnings and the exact selection procedure.
Named-edge conversion is preferable to guessing direction from an unlabeled
binary matrix. An incomplete or nonextendible graph must remain an explicit
scientific or computational limitation, with no fallback to a conveniently
oriented causal story.

\subsection{The permitted prototype path and its failure contract}

The proposed operation accepts one exposure and one outcome in a fully
observed acyclic graph, with unique string node identifiers and an explicit
node list. It initially caps the graph at twelve nodes as a conservative
wrapper design choice; that number is not measured engine capacity.
Latent, bidirected, selection, undirected and ambiguous-edge inputs are
outside this path. It constructs a fresh pinned DAG with explicit roles
and requests the minimal adjustment candidate. The original graph remains
available for comparison \cite{exec13_pgmpy_reuse_oracle}.

After a successful return, the wrapper must verify unchanged nodes, edges,
exposure and outcome roles; validate candidate membership and timing;
exclude exposure descendants; and run the independent graphical check.
Success then means a candidate satisfying the specified criterion relative
to the supplied graph. It carries no numerical effect, confidence interval
or certification of the graph's real-world truth. An empty valid set is
a meaningful result, whereas a null or failed return is a different
disposition. A failure to find a set by this method must not become a
claim of global causal nonidentifiability
\cite{exec13_pgmpy_reuse_oracle}.

Exceptions, warnings, unsupported variants, altered graph structure and
resource limits are retained with their causes. They do not become an
empty adjustment set or zero effect. The same graph, target identity and
disposition must cross library, CLI and MCP boundaries. Implementation
acceptance will require actual pinned-engine calls on named confounder,
mediator, collider, isolate and permutation fixtures, together with invalid
graph and failure cases. The exact references developed here provide
oracles for that work; no pgmpy runtime or adapter call was performed
\cite{exec13_pgmpy_reuse_oracle}.

The minimal path's attraction is limited but practical: pgmpy can propose
a graph-based candidate while the project's independent sufficient check
guards its meaning. It is not necessary to copy the whole library or
port its graph stack to another language to obtain this benefit. A
reference-only role remains appropriate for presently unverified discovery
and estimation branches. Minimum-variance adjustment is rejected at this
pin because the branch is unimplemented. Generalized adjustment, latent
graphs and data-driven direction selection are deferred to separate
source and semantic reviews \cite{exec13_pgmpy_reuse_primary}.

\subsection{Licensing, dependencies and ownership make reuse concrete}

The pinned MIT license permits use, modification and distribution subject
to preservation of its copyright and permission notice and contains a
warranty disclaimer. That permission does not replace dependency-license
review. Project metadata requires Python from 3.10 up to but excluding
3.15 and declares twelve mandatory dependencies, including NumPy,
SciPy, NetworkX, pandas, scikit-learn, statsmodels, scikit-base and
Hugging Face Hub. Optional groups separately include Torch, graphical
tools and LiteLLM. Declared version floors are not a resolved reproducible
environment \cite{exec13_pgmpy_reuse_primary}.

The proposed graph service should use an isolated, locked environment and
only the necessary path. Owners must retain source and dependency identities,
notices, conversion fixtures and resource measurements. An upgrade requires
rechecking API names, graph representations and failure behavior as well as
rerunning accepted fixtures. The inspected revision includes adjustment
tests for observed, latent and M-graph cases, which are useful maintenance
evidence; their presence does not imply that those tests were executed here
or that all targets and graph classes are qualified
\cite{exec13_pgmpy_reuse_primary}.

The LLM pairwise class is also a separate capability within the same
repository. Its description states that it uses variable names and text
descriptions, not the observed data values, and its source includes an
optional provider-call path. Such suggestions can inform a reviewable
business formulation but do not inherit the evidence of a CI-based
discovery algorithm. This selected graph operation has no live-model
dependency. A future language-model path would need its own custody,
ambiguity, source-grounding and held-out evaluation contract
\cite{exec13_pgmpy_reuse_primary}.

Costs therefore include graph elicitation, controls and timing review,
dependency maintenance, integration fixtures and adjudication of failures,
even though the selected source license has no purchase price. No monetary
maintenance estimate is supplied without workload evidence. Research owners
must justify the assumed graph and target; adapter owners must enforce
the conversion and failure contract; reviewers must verify executed
results before runtime qualification. The concrete revisit trigger is
a new pin, graph class or method path, followed by source inspection and
matched behavioral tests. This division gives pgmpy a useful, limited
place in the proposal while keeping the unimplemented and scientifically
unsupported extensions visible.
\label{exec13:pgmpy-reuse-end}
